% Part: second-order-logic
% Chapter: metatheory
% Section: compactness

\documentclass[../../../include/open-logic-section]{subfiles}

\begin{document}

\olfileid[pt]{sol}{met}{com}

\olsection{A lógica de segunda ordem não é compacta}

\begin{explain}
Chamamos um conjunto de sentenças~$\Gamma$ de \emph{finitamente
satisfatível} se cada um de seus subconjuntos finitos é satisfatível. A
lógica de primeira ordem tem a propriedade de que, se um conjunto de
!!{sentence}s~$\Gamma$ é finitamente satisfatível, então é satisfatível.
Essa propriedade se chama \emph{compacidade}. Ela tem uma versão equivalente
que envolve consequência: se $\Gamma \Entails !A$, então já temos
$\Gamma_0 \Entails !A$ para algum subconjunto finito
~$\Gamma_0 \subseteq \Gamma$. Nessa versão, ela é um corolário imediato do
teorema da completude: pois, se $\Gamma \Entails !A$, pela completude
$\Gamma \Proves !A$. Mas !!{derivation} só pode usar um número finito de
!!{sentence}s de~$\Gamma$.

A compacidade não vale para a lógica de segunda ordem. Há conjuntos de
!!{sentence}s de segunda ordem que são finitamente satisfatíveis, mas não
satisfatíveis, e que acarretam algum~$!A$ sem que um subconjunto finito
acarrete~$!A$.
\end{explain}


\begin{thm}
\ollabel{thm:sol-undecidable}
A lógica de segunda ordem não é compacta.
\end{thm}

\begin{proof}
Lembre-se de que
\[
\fn{Inf} \ident \lexists[u][(\lforall[x][\lforall[y][(\eq[u(x)][u(y)]
      \lif \eq[x][y])]] \land \lexists[y][\lforall[x][\eq/[y][u(x)]]])]
\]
é satisfeita em !!a{structure} sse seu domínio é infinito. Seja
$!A^{\ge n}$ uma sentença que afirma que o domínio tem ao menos
$n$ !!{element}s, por exemplo,
\[
!A^{\ge n} \ident \lexists[x_1][\dots\lexists[x_n][(\eq/[x_1][x_2]
    \land \eq/[x_1][x_3] \land \dots \land \eq/[x_{n-1}][x_n])]].
\]
Considere o conjunto de !!{sentence}s
\[
\Gamma = \{\lnot \fn{Inf}, !A^{\ge 1}, !A^{\ge 2}, !A^{\ge 3}, \dots\}.
\]
Ele é finitamente satisfatível, pois, para qualquer subconjunto finito
~$\Gamma_0 \subseteq \Gamma$, existe algum $k$ tal que
$!A^{\ge k} \in \Gamma$, mas nenhum $!A^{\ge n} \in \Gamma$ para $n>k$.
Se $\Domain{M}$ tem $k$ !!{element}s, então $\Sat{M}{\Gamma_0}$. Mas
$\Gamma$ não é satisfatível: se $\Sat{M}{\lnot \fn{Inf}}$, então
$\Domain{M}$ deve ser finito, digamos, de tamanho~$k$. Nesse caso,
$\Sat/{M}{!A^{\ge k+1}}$.
\end{proof}

\begin{prob}
Dê um exemplo de um conjunto~$\Gamma$ e de !!a{sentence}~$!A$ tais que
$\Gamma \Entails !A$, mas, para todo subconjunto finito
~$\Gamma_0 \subseteq \Gamma$, $\Gamma_0 \Entails/ !A$.
\end{prob}


\end{document}
